\section{Introduction}\label{sec:intro}
Power-system operators release forecasts, prices, reserve requirements, and operating statistics to support transparent and efficient operation. These records describe a shared physical system, so their information content is interdependent. A price series and a renewable forecast may each reveal little about a generating unit; together, they can help reconstruct the residual demand that determines its dispatch. A release decision made separately for each field can therefore overlook the information exposed by their combination.

This problem is particularly relevant to data classification and grading. Disclosure rules specify which operational fields are released and when~\cite{nea2024disclosure}, while classification standards call for aggregation and processing effects to be considered~\cite{gbt43697}. The audit nevertheless needs a usable output at the field level: which field poses a risk, how much risk it adds, and which other fields make that risk material. The inference and aggregation problem has long been recognized in database security~\cite{lunt1989aggregation,farkas2002inference}; in power systems, even public prices can reveal grid topology~\cite{kekatos2014topology}. Translating this set-dependent leakage into an actionable field-level assessment remains the central challenge.

Existing scores address parts of the problem. Correlation and single-field prediction measure what one field reveals alone. Feature-attribution methods distribute predictive value or evaluate a field in a particular fitted model. Inference-graph methods propagate sensitivity through relations between fields~\cite{hu2026ignn}. For a release audit, however, the relevant question is whether a field can substantially improve an adversary's inference in \emph{any admissible background}. A field may have low individual predictability, little average attribution, and yet be decisive in one combination.

We take the predictive accuracy of attackers trained on a field set as the starting point. A field's \emph{budgeted marginal inference risk} is its largest increase in this accuracy over backgrounds containing at most $K$ other fields, relative to the public baseline. The score adapts the marginal contribution index~\cite{catav2021mci} to a bounded-background release audit. Its maximizing background supplies an explanation: the risk score is accompanied by the fields with which the contribution arises. Threshold-crossing combinations then connect the assessment to field flagging and withholding.

The computational difficulty is evaluating many such backgrounds. Even 35 candidate fields produce 7,175 nonempty sets of size at most three. We therefore share a nonlinear representation across sets and fit a closed-form readout separately for each query. This amortized attacker scans the field combinations; dedicated retraining of a few selected backgrounds verifies the largest contributions. Fig.~\ref{fig:framework} summarizes the resulting audit.

The contributions are threefold:
\begin{itemize}
\item A field-level risk measure that retains the combination structure of set-level inferability. We establish its safety-margin interpretation and its connection to minimal unsafe sets under monotonicity, and use witness backgrounds to explain field risk.
\item An amortized estimator based on masked pre-training and subset-specific readouts, together with a scan--certify procedure that recovers retrained lower bounds from a small number of backgrounds.
\item An evaluation over four power-system data sets, ten targets, and 124,425 field sets. The experiments measure missed combination risk, compare grading and withholding strategies, isolate the estimator's effective components, and assess robustness and background-budget coverage.
\end{itemize}

\begin{figure*}[!tp]
\centering
\includegraphics[width=\textwidth]{fig1_framework.pdf}
\caption{From operational data to a field-level audit. The grid and traces illustrate the shared physical origin of the released fields. A field is evaluated in combination with a background $T$ relative to the public baseline $B$; the largest gain defines its risk and supplies a witness. An amortized scan selects backgrounds for retrained assessment, while unsafe combinations support withholding decisions. The grid and combinations are schematic.}
\label{fig:framework}
\end{figure*}
